% GENERATED FILE. Edit canonical lessons, then run npm run generate:books.
% canonical-source-hash: fnv1a64:0bb0bc7925541392
% canonical-lessons: FA-C521-sher, FA-C521-shaer, FA-C521-ghazal, FA-C521-robai, FA-C521-beyt

\chapter{Poetry}
\label{ch:fa-poetry}
\hlchaptermodality{521}

\begin{chapteropening}
\textbf{By the end of this chapter:} \emph{I can talk about poems and poets.}

Complete the last lesson of chapter 521: I can talk about poems and poets.
\end{chapteropening}

\section[\texorpdfstring{\fa{شعر} (she'r)}{شعر (she'r)}]{\fa{شعر} (she'r) — a poem; poetry}

\label{lesson:FA-C521-sher}



\begin{quote}
Before the new one: say the Persian for a choir, then the Persian for an album.
\end{quote}

\subsection*{You'll want to know: \fa{شعر}}
\textbf{\fa{شعر}} (\emph{she'r}) — \textquotedblleft{}a poem; poetry\textquotedblright{}.

\begin{quote}
\fa{او} \fa{هر} \fa{شب} \fa{یک} \fa{شعر} \fa{میخواند}.
\end{quote}

\begin{quote}
\emph{u har shab yek she'r mikhânad.}
\end{quote}

\begin{quote}
He reads a poem every night.
\end{quote}

Plural \emph{she'rhâ}, or the Arabic \emph{ash'âr}. \emph{she'r goftan} — to write poetry.

\begin{cousinweb}
\emph{she'r} comes from Arabic \emph{šiʿr}, 'poem, poetry' (EtymDB).
\end{cousinweb}

\begin{grammarlens}[title={What you've built}]
One more word for telling, explaining and describing at B1.
\end{grammarlens}

\subsection*{Your turn}
\begin{itemize}
\raggedright
  \item \emph{Say it:} \emph{she'r}
  \item \emph{Say it:} \emph{she'r}, once more
  \item \emph{Say it:} a sentence of your own with \emph{she'r}
  \item \emph{Recall:} say the Persian for a choir, then the Persian for an album, then say \emph{she'r} again
\end{itemize}

\begin{tcolorbox}[breakable,colback=teal!4,colframe=teal!35!black,title={Before you move on}]
What does \fa{شعر} mean? (\textquotedblleft{}A poem; poetry\textquotedblright{}.) Say it once more.
\end{tcolorbox}
\section[\texorpdfstring{\fa{شاعر} (shâ'er)}{شاعر (shâ'er)}]{\fa{شاعر} (shâ'er) — a poet}

\label{lesson:FA-C521-shaer}



\begin{quote}
Before the new one: say the Persian for an album, then the Persian for a poem.
\end{quote}

\subsection*{You'll want to know: \fa{شاعر}}
\textbf{\fa{شاعر}} (\emph{shâ'er}) — \textquotedblleft{}a poet\textquotedblright{}.

\begin{quote}
\fa{پدرم} \fa{شاعر} \fa{است} \fa{و} \fa{کتاب} \fa{شعر} \fa{دارد}.
\end{quote}

\begin{quote}
\emph{pedaram shâ'er ast va ketâb-e she'r dârad.}
\end{quote}

\begin{quote}
My father is a poet and has a book of poems.
\end{quote}

Plural \emph{shâ'erân}, or the Arabic \emph{sho'arâ}.

\begin{cousinweb}
\emph{shâ'er} was borrowed from Arabic \emph{šāʿir}, 'poet' (EtymDB), a word from the same root as \emph{she'r}.
\end{cousinweb}

\begin{grammarlens}[title={What you've built}]
One more word for telling, explaining and describing at B1.
\end{grammarlens}

\subsection*{Your turn}
\begin{itemize}
\raggedright
  \item \emph{Say it:} \emph{shâ'er}
  \item \emph{Say it:} \emph{shâ'er}, once more
  \item \emph{Say it:} a sentence of your own with \emph{shâ'er}
  \item \emph{Recall:} say the Persian for an album, then the Persian for a poem, then say \emph{shâ'er} again
\end{itemize}

\begin{tcolorbox}[breakable,colback=teal!4,colframe=teal!35!black,title={Before you move on}]
What does \fa{شاعر} mean? (\textquotedblleft{}A poet\textquotedblright{}.) Say it once more.
\end{tcolorbox}
\section[\texorpdfstring{\fa{غزل} (ghazal)}{غزل (ghazal)}]{\fa{غزل} (ghazal) — a ghazal (a love lyric)}

\label{lesson:FA-C521-ghazal}



\begin{quote}
Before the new one: say the Persian for a poem, then the Persian for a poet.
\end{quote}

\subsection*{You'll want to know: \fa{غزل}}
\textbf{\fa{غزل}} (\emph{ghazal}) — \textquotedblleft{}a ghazal (a love lyric)\textquotedblright{}.

\begin{quote}
\fa{مادربزرگم} \fa{غزلهای} \fa{زیادی} \fa{را} \fa{از} \fa{حفظ} \fa{میداند}.
\end{quote}

\begin{quote}
\emph{mâdarbozorgam ghazalhâ-ye ziyâdi râ az hefz midânad.}
\end{quote}

\begin{quote}
My grandmother knows many ghazals by heart.
\end{quote}

A short lyric poem in couplets, usually about love. Plural \emph{ghazalhâ}.

\begin{cousinweb}
\emph{ghazal} was borrowed from Arabic \emph{ġazal}, which EtymDB links with a verb meaning to speak words of love to the beloved (EtymDB). English \emph{ghazal} is usually traced to Persian and Urdu.
\end{cousinweb}

\begin{grammarlens}[title={What you've built}]
One more word for telling, explaining and describing at B1.
\end{grammarlens}

\subsection*{Your turn}
\begin{itemize}
\raggedright
  \item \emph{Say it:} \emph{ghazal}
  \item \emph{Say it:} \emph{ghazal}, once more
  \item \emph{Say it:} a sentence of your own with \emph{ghazal}
  \item \emph{Recall:} say the Persian for a poem, then the Persian for a poet, then say \emph{ghazal} again
\end{itemize}

\begin{tcolorbox}[breakable,colback=teal!4,colframe=teal!35!black,title={Before you move on}]
What does \fa{غزل} mean? (\textquotedblleft{}A ghazal (a love lyric)\textquotedblright{}.) Say it once more.
\end{tcolorbox}
\section[\texorpdfstring{\fa{رباعی} (robâ'i)}{رباعی (robâ'i)}]{\fa{رباعی} (robâ'i) — a quatrain (four-line poem)}

\label{lesson:FA-C521-robai}



\begin{quote}
Before the new one: say the Persian for a poet, then the Persian for a ghazal.
\end{quote}

\subsection*{You'll want to know: \fa{رباعی}}
\textbf{\fa{رباعی}} (\emph{robâ'i}) — \textquotedblleft{}a quatrain (four-line poem)\textquotedblright{}.

\begin{quote}
\fa{او} \fa{یک} \fa{رباعی} \fa{کوتاه} \fa{برای} \fa{مادرش} \fa{نوشت}.
\end{quote}

\begin{quote}
\emph{u yek robâ'i-ye kutâh barâ-ye mâdarash nevesht.}
\end{quote}

\begin{quote}
He wrote a short quatrain for his mother.
\end{quote}

Four lines, the third free of the rhyme. Plural \emph{robâ'iyât}.

\begin{cousinweb}
\emph{robâ'i} is a loan, usually traced to Arabic \emph{rubāʿī}, 'fourfold', from the Arabic word for four. EtymDB has no entry for it.
\end{cousinweb}

\begin{grammarlens}[title={What you've built}]
One more word for telling, explaining and describing at B1.
\end{grammarlens}

\subsection*{Your turn}
\begin{itemize}
\raggedright
  \item \emph{Say it:} \emph{robâ'i}
  \item \emph{Say it:} \emph{robâ'i}, once more
  \item \emph{Say it:} a sentence of your own with \emph{robâ'i}
  \item \emph{Recall:} say the Persian for a poet, then the Persian for a ghazal, then say \emph{robâ'i} again
\end{itemize}

\begin{tcolorbox}[breakable,colback=teal!4,colframe=teal!35!black,title={Before you move on}]
What does \fa{رباعی} mean? (\textquotedblleft{}A quatrain (four-line poem)\textquotedblright{}.) Say it once more.
\end{tcolorbox}
\section[\texorpdfstring{\fa{بیت} (beyt)}{بیت (beyt)}]{\fa{بیت} (beyt) — a couplet, a line of verse}

\label{lesson:FA-C521-beyt}



\begin{quote}
Before the new one: say the Persian for a ghazal, then the Persian for a quatrain.
\end{quote}

\subsection*{You'll want to know: \fa{بیت}}
\textbf{\fa{بیت}} (\emph{beyt}) — \textquotedblleft{}a couplet, a line of verse\textquotedblright{}.

\begin{quote}
\fa{او} \fa{یک} \fa{بیت} \fa{از} \fa{این} \fa{غزل} \fa{را} \fa{خواند}.
\end{quote}

\begin{quote}
\emph{u yek beyt az in ghazal râ khând.}
\end{quote}

\begin{quote}
He recited a couplet from this ghazal.
\end{quote}

Plural \emph{abyât} (Arabic). A \emph{beyt} is two half-lines.

\begin{cousinweb}
\emph{beyt} was borrowed from Arabic \emph{bayt}, 'house' (EtymDB), which EtymDB traces on to Proto-Semitic \emph{bayt-}. In poetry the couplet is a little house of words.
\end{cousinweb}

\begin{grammarlens}[title={What you've built}]
One more word for telling, explaining and describing at B1.
\end{grammarlens}

\subsection*{Your turn}
\begin{itemize}
\raggedright
  \item \emph{Say it:} \emph{beyt}
  \item \emph{Say it:} \emph{beyt}, once more
  \item \emph{Say it:} a sentence of your own with \emph{beyt}
  \item \emph{Recall:} say the Persian for a ghazal, then the Persian for a quatrain, then say \emph{beyt} again
\end{itemize}

\begin{tcolorbox}[breakable,colback=teal!4,colframe=teal!35!black,title={Before you move on}]
What does \fa{بیت} mean? (\textquotedblleft{}A couplet, a line of verse\textquotedblright{}.) Say it once more.
\end{tcolorbox}
